% Rekonstruktion von Funktionsgleichungen – bank/rekonstruktion-von-funktionsgleichungen/ (zusammenbau v0.4, Heft)
\einheitenkopf[t4]{Rekonstruktion von Funktionsgleichungen}
\setcounter{aufgabe}{13}

% Anlauf (ohne Prüfkennung)
\begin{aufgabe}{Ansatz und Punktbedingungen – Anlauf}
\begin{teile}
\teil Gesucht ist eine quadratische Funktion $f(x) = ax^2 + bx + c$; bekannt sind drei Punkte des Graphen. Reichen die Angaben? \\ \kreuz{zu wenig} \\ \kreuz{genau passend} \\ \kreuz{mehr als nötig}
\teil Bestimme die Gleichung der Geraden durch $A(1|5)$ und $B(4|14)$. y = \leerfeld
\teil Der Graph einer ganzrationalen Funktion $f$ vierten Grades ist achsensymmetrisch zur $y$-Achse und verläuft durch $P(0|4)$, $Q(1|1)$ und $R(2|4)$. Bestimme eine Funktionsgleichung von $f$. f(x) = \leerfeld
\end{teile}
\end{aufgabe}

\begin{aufgabe}{Ansatz und Punktbedingungen – Prüfungsaufgaben}
\begin{teile}
\teil Die Gerade $h$ verläuft durch $C(0{,}6|1{,}2)$ und $D(3|6)$. Weise nach, dass $h(x) = 2x$ eine Gleichung von $h$ ist. \hfill (Abitur 2022 GK)
\teil Eine Seilbahn verläuft im Modell geradlinig durch die Stützen $E(-1|2{,}5)$ und $F(3|0{,}5)$. Weise nach, dass $k(x) = -0{,}5x + 2$ die Gleichung dieser Geraden ist. \hfill (Abitur 2022 GK)
\teil Die Gerade $s$ verläuft durch $G(1{,}5|-1)$ und $H(3{,}5|3)$. Weise nach, dass $s(x) = 2x - 4$ eine Gleichung von $s$ ist. \hfill (Abitur 2022 GK)
\teil Die Abbildung zeigt den Graphen von $f$ mit $f(x) = a \cdot b^x$ und $a, b > 0$. Lies zwei geeignete Punkte ab und bestimme $a$ und $b$. \hfill (Abitur 2022 GK) \\

\begin{ksys}[xmin=-3,xmax=2,ymin=0,ymax=7,ablesen] \funktionab{2*3^\x}{f}{-3}{1.1} \end{ksys}
a = \leerfeld , b = \leerfeld
\end{teile}
\end{aufgabe}

% Anlauf (ohne Prüfkennung)
\begin{aufgabe}{Bedingungen mit Ableitung – Anlauf}
\begin{teile}
\teil „Der Graph von $f$ hat den Hochpunkt $H(3|1)$.“ Wie viele Gleichungen stecken in diesem Satz? \\ \kreuz{eine Gleichung} \\ \kreuz{zwei Gleichungen}
\teil Der Graph einer quadratischen Funktion $f$ verläuft durch $A(0|-1)$ und hat den Tiefpunkt $T(1|-3)$. Bestimme eine Funktionsgleichung von $f$. f(x) = \leerfeld
\teil Der Graph einer quadratischen Funktion $f$ schneidet die $y$-Achse bei $y = 4$; die Tangente an den Graphen an der Stelle $x = -1$ hat die Gleichung $y = 4x + 7$. Bestimme einen Funktionsterm von $f$. f(x) = \leerfeld
\end{teile}
\end{aufgabe}

\begin{aufgabe}{Bedingungen mit Ableitung – Prüfungsaufgaben}
\begin{teile}
\teil Die untere Begrenzung eines Blumenbeets soll durch eine Parabel $q$ beschrieben werden, deren Graph durch den Ursprung verläuft und den Hochpunkt $H(3|2{,}7)$ hat. Bestimme eine Funktionsgleichung von $q$. \hfill (Abitur 2021 GK) \\ q(x) = \leerfeld
\teil Die Höhe einer Sonnenblume wird in den ersten Wochen durch $h(t) = a \cdot t^3 + b \cdot t^2$ beschrieben ($t$ in Wochen, $h(t)$ in cm). Nach 4 Wochen ist sie 64 cm hoch und wächst in diesem Moment mit 24 cm pro Woche. Bestimme eine Funktionsgleichung von $h$. \hfill (Abitur 2019 GK) \\ h(t) = \leerfeld
\teil Die Temperatur in einem Gewächshaus am Vormittag wird durch eine ganzrationale Funktion $w$ dritten Grades beschrieben ($x$ in Stunden, $w(x)$ in °C) mit: I $w(0) = 15$, II $w'(0) = 4$, III $H(4|31)$ ist Hochpunkt des Graphen von $w$. Bestimme eine Funktionsgleichung von $w$. \hfill (Abitur 2023 GK) \\ w(x) = \leerfeld
\teil Der Graph einer quadratischen Funktion $f$ verläuft durch den Ursprung; die Tangente an den Graphen im Punkt $(1|f(1))$ hat die Gleichung $y = 3x + 1$. Bestimme einen Funktionsterm von $f$. \hfill (Abitur 2020 GK) \\ f(x) = \leerfeld
\teil Der Graph einer quadratischen Funktion $f$ verläuft durch den Ursprung; die Tangente an den Graphen im Punkt $(3|f(3))$ hat die Gleichung $y = -2x + 9$. Bestimme einen Funktionsterm von $f$. \hfill (Abitur 2020 GK) \\ f(x) = \leerfeld
\teil Bei einer Dirtbike-Schanze endet die Absprungrampe $r$ mit $r(x) = 0{,}1x^2 + 2{,}5$ ($x \le 0$, alle Angaben in m) im Punkt $S(0|2{,}5)$; dort geht sie ohne Knick in die Flugbahn über, die durch eine quadratische Funktion $f$ beschrieben wird. Der Landehügel wird durch $g(x) = 2{,}5 - 0{,}5x$ beschrieben. Bei $x = 4$ ist der Fahrer 1,2 m senkrecht über dem Landehügel. Bestimme eine Funktionsgleichung der Flugbahn $f$. \hfill (Abitur 2018 GK) \\ f(x) = \leerfeld
\teil Die Kante eines Blechteils verläuft vom Ursprung bis $R(1|3)$ gekrümmt und ab $R$ geradlinig mit der Steigung 8. Das gekrümmte Stück soll durch $k(x) = a \cdot x^4 + b \cdot x^2$ beschrieben werden und in $R$ ohne Knick in die gerade Kante übergehen. Bestimme $a$ und $b$. \hfill (Abitur 2022 GK) \\ a = \leerfeld , b = \leerfeld
\end{teile}
\end{aufgabe}

\begin{aufgabe}{Zwei Parameter einer Exponentialfunktion aus Funktionswert und Änderungsrate an einer Stelle bestimmen – Prüfungsaufgaben}
\begin{teile}
\teil Eine Suppe wird vom Herd genommen und kühlt ab. Ab $t = 5$ ($t$ in Minuten) wird ihre Temperatur durch $w(t) = 20 + b \cdot e^{c \cdot t}$ beschrieben ($w(t)$ in °C). Zu Beginn des Abkühlens beträgt die Temperatur 80 °C und die momentane Änderungsrate $-6$ °C pro Minute. Bestimme $b$ und $c$; runde $b$ auf zwei Stellen. \hfill (Abitur 2017 GK) \\ b ≈ \leerfeld , c = \leerfeld
\teil Ein Joghurt wird aus dem Kühlschrank genommen und erwärmt sich. Ab $t = 2$ ($t$ in Stunden) wird seine Temperatur durch $w(t) = 22 + b \cdot e^{c \cdot t}$ beschrieben ($w(t)$ in °C). Zu Beginn des Erwärmens beträgt die Temperatur 7 °C und die momentane Änderungsrate 3 °C pro Stunde. Bestimme $b$ und $c$; runde $b$ auf zwei Stellen. \hfill (Abitur 2017 GK) \\ b ≈ \leerfeld , c = \leerfeld
\teil Ein Gartenteich kühlt am Abend ab. Ab $t = 4$ ($t$ in Stunden seit 16 Uhr) wird seine Wassertemperatur durch $w(t) = 12 + b \cdot e^{c \cdot t}$ beschrieben ($w(t)$ in °C). Zu Beginn des Abkühlens beträgt sie 18 °C, die momentane Änderungsrate ist $-1{,}5$ °C pro Stunde. Bestimme $b$ und $c$; runde $b$ auf zwei Stellen. \hfill (Abitur 2017 GK) \\ b ≈ \leerfeld , c = \leerfeld
\end{teile}
\end{aufgabe}

% Anlauf (ohne Prüfkennung)
\begin{aufgabe}{Sonderansätze und Modellkritik – Anlauf}
\begin{teile}
\teil Ein Sinusgraph hat einen Hochpunkt bei $x = 1$ und den nächsten Tiefpunkt bei $x = 4$. Wie weit ist es vom Hochpunkt bis zum nächsten Hochpunkt? \\ \kreuz{$3$} \\ \kreuz{$6$} \\ \kreuz{$12$}
\teil Der Graph von $f$ mit $f(x) = a \cdot \mathrm{sin}(b \cdot x)$ und $a, b > 0$ hat die direkt aufeinanderfolgenden Extrempunkte $E_1(1|3)$ und $E_2(3|-3)$. Bestimme $a$ und $b$. a = \leerfeld , b = \leerfeld
\teil Der Wasserstand an einem Pegel, gemessen gegenüber dem Mittelwasser, wird durch $w(t) = a \cdot \mathrm{sin}(b \cdot t)$ mit $a, b > 0$ beschrieben ($t$ in Stunden, $w(t)$ in m). Zum Zeitpunkt $t = 0$ ist Mittelwasser, das erste Hochwasser folgt nach 3,1 Stunden und liegt 1,8 m über dem Mittelwasser. Bestimme $a$ und $b$ auf zwei Stellen. a = \leerfeld , b ≈ \leerfeld
\end{teile}
\end{aufgabe}

\begin{aufgabe}{Sonderansätze und Modellkritik – Prüfungsaufgaben}
\begin{teile}
\teil Die Funktion $f$ mit $f(x) = x - \frac{1}{27}x^3$ soll auf $[-3; 3]$ durch $g$ mit $g(x) = a \cdot \mathrm{sin}(b \cdot x)$ und $a, b > 0$ angenähert werden. Dabei ist $3$ die erste positive Extremstelle von $g$, und es gilt $g(3) = f(3)$. Bestimme $a$ und $b$. \hfill (Abitur 2022 GK) \\ a = \leerfeld , b = \leerfeld
\teil Die Funktion $f$ mit $f(x) = 1{,}5x - \frac{1}{32}x^3$ soll auf $[-4; 4]$ durch $g$ mit $g(x) = a \cdot \mathrm{sin}(b \cdot x)$ und $a, b > 0$ angenähert werden. Dabei ist $4$ die erste positive Extremstelle von $g$, und es gilt $g(4) = f(4)$. Bestimme $a$ und $b$. \hfill (Abitur 2022 GK) \\ a = \leerfeld , b = \leerfeld
\end{teile}
\end{aufgabe}

\begin{aufgabe}{Existenz einer quadratischen Funktion zu vier Wert- und Steigungsbedingungen über das überbestimmte Gleichungssystem untersuchen – Prüfungsaufgaben}
\begin{teile}
\teil Der Graph einer Funktion $f$ verläuft durch $A(0|3)$ und $B(1|0{,}5)$ mit $f'(0) = -2$ und $f'(1) = 1$. Der Graph einer quadratischen Funktion $p$ soll in $A$ und in $B$ tangential zum Graphen von $f$ verlaufen. Stelle vier Bedingungen für $p$ auf und untersuche, ob es eine solche Funktion $p$ gibt. \hfill (Abitur 2017 GK)
\teil Der Graph einer Funktion $f$ verläuft durch $A(0|-1)$ und $B(1|1)$ mit $f'(0) = 1$ und $f'(1) = 3$. Der Graph einer quadratischen Funktion $p$ soll in $A$ und in $B$ tangential zum Graphen von $f$ verlaufen. Stelle vier Bedingungen für $p$ auf und untersuche, ob es eine solche Funktion $p$ gibt. \hfill (Abitur 2017 GK)
\teil Der Graph einer Funktion $f$ verläuft durch $A(0|4)$ und $B(2|0)$ mit $f'(0) = 0$ und $f'(2) = -3$. Der Graph einer quadratischen Funktion $p$ soll in $A$ und in $B$ tangential zum Graphen von $f$ verlaufen. Stelle vier Bedingungen für $p$ auf und untersuche, ob es eine solche Funktion $p$ gibt. \hfill (Abitur 2017 GK)
\end{teile}
\end{aufgabe}
